% 第 1 章：集合与常用逻辑用语 · 笔记
% 笔记写法约定见 chapters/README.md
\section{集合与常用逻辑用语}

\begin{notezone}{课前：我已知什么}
\begin{itemize}
  \item 初中见过的数集：\underline{\hspace{4em}}、\underline{\hspace{4em}}、\underline{\hspace{4em}}；
  \item 「所有大于 3 的数」这种说法，能不能写成一个集合？\hspace{2em}（能 / 不能）
\end{itemize}
\noteblank[10ex]
\end{notezone}

\begin{notezone}{课中：概念怎么说的}
\begin{itemize}
  \item 元素的三个特性：\hspace{8em}；
  \item 列举法与描述法的区别：\hspace{8em}；
  \item 描述法 $\{x\mid p(x)\}$ 里，$x$ 叫\hspace{4em}，$p(x)$ 叫\hspace{4em}。
\end{itemize}
\noteblank[14ex]
\end{notezone}

\begin{notezone}{课中：一道我没听懂的例题}
（抄题或写题号）

\noteblank[16ex]

我当时卡在哪一步？

\noteblank[10ex]
\end{notezone}

\begin{notezone}{课后：我能复述什么}
试着不看书画出「集合关系 ↔ 条件关系」的对应表：

\noteblank[18ex]

今天最容易错的一点是：

\noteblank[8ex]
\end{notezone}
